\section{Confidence-Gated Self-Training}
\label{sec:method}

\subsection{Problem setting}

Let $\mathcal{D}_s=\{(x_i^s,y_i^s)\}_{i=1}^{n_s}$ be labelled source data and
$\mathcal{D}_t=\{x_j^t\}_{j=1}^{n_t}$ unlabelled target data over $K$ classes.
A classifier $f_\theta:\mathcal{X}\rightarrow\Delta^{K-1}$ outputs a categorical
distribution. We minimise the usual supervised source loss
\begin{equation}
  \mathcal{L}_{s}(\theta)
  = \E_{(x,y)\sim\mathcal{D}_s}
    \bigl[-\log f_\theta(x)_y\bigr].
  \label{eq:source}
\end{equation}
For a target example, a weak augmentation $a_w$ produces
$p=f_\theta(a_w(x))$. Its candidate pseudo-label and confidence are
$\hat y=\argmax_k p_k$ and $q=\max_k p_k$, respectively.

\subsection{Adaptive confidence gate}

At training step $r$, we maintain an exponential moving average $m_k^{(r)}$ of
the confidence assigned to examples predicted as class $k$:
\begin{equation}
  m_k^{(r)}=\rho m_k^{(r-1)}+(1-\rho)
  \frac{\sum_{x\in B_t}\1[\hat y(x)=k]q(x)}
       {\max\{1,\sum_{x\in B_t}\1[\hat y(x)=k]\}},
  \label{eq:ema}
\end{equation}
with momentum $\rho$. Empty classes retain their preceding value. The
class-conditional acceptance threshold is
\begin{equation}
  \tau_k^{(r)}=\operatorname{clip}
  \left(\tau_0\frac{m_k^{(r)}}{\frac{1}{K}\sum_{j=1}^{K}m_j^{(r)}},
  \tau_{\min},\tau_{\max}\right).
  \label{eq:threshold}
\end{equation}
The normalization makes thresholds lower for classes that are systematically
less confident, while clipping prevents degenerate acceptance.

\subsection{Augmentation agreement and objective}

Let $p'=f_\theta(a_s(x))$ denote the prediction under a strong augmentation.
The reliability gate is
\begin{equation}
  g(x)=\1[q(x)\ge\tau_{\hat y(x)}]
       \1[\argmax_k p'_k=\hat y(x)].
  \label{eq:gate}
\end{equation}
We stop gradients through $p$, $\hat y$, and $g$. The target loss is the
coverage-normalised cross-entropy
\begin{equation}
  \mathcal{L}_t(\theta)=
  \frac{\sum_{x\in B_t}g(x)\bigl[-\log p'_{\hat y(x)}\bigr]}
       {\max\{1,\sum_{x\in B_t}g(x)\}}.
  \label{eq:target}
\end{equation}
The final objective is
$\mathcal{L}=\mathcal{L}_s+\lambda(r)\mathcal{L}_t$, where $\lambda(r)$ is
linearly warmed from zero to one during the first $10\%$ of training. The gate
is used only during training, so inference cost equals that of the backbone.

\begin{figure}[t]
  \centering
  \begin{tikzpicture}[
      node distance=5mm and 8mm,
      box/.style={draw=linkblue,rounded corners=2pt,fill=softblue,
                  align=center,minimum height=8mm,font=\small},
      decision/.style={draw=darkgray,rounded corners=2pt,align=center,
                       minimum height=8mm,font=\small},
      arrow/.style={-{Latex[length=2mm]},thick,draw=darkgray}
    ]
    \node[box] (target) {unlabelled\\target image $x$};
    \node[box,right=of target] (weak) {weak view\\$a_w(x)$};
    \node[box,below=of weak] (strong) {strong view\\$a_s(x)$};
    \node[decision,right=of weak] (conf) {class-adaptive\\confidence test};
    \node[decision,right=of strong] (agree) {prediction\\agreement test};
    \node[box,right=11mm of $(conf)!0.5!(agree)$] (loss) {accepted\\target loss};
    \draw[arrow] (target) -- (weak);
    \draw[arrow] (target.east) -- ++(3mm,0) |- (strong.west);
    \draw[arrow] (weak) -- (conf);
    \draw[arrow] (strong) -- (agree);
    \draw[arrow] (conf.east) -- ++(3mm,0) |- (loss.west);
    \draw[arrow] (agree.east) -- ++(3mm,0) |- (loss.west);
  \end{tikzpicture}
  \caption{Overview of \method. A pseudo-label contributes to the target loss
  only when it passes both confidence and augmentation-agreement tests.}
  \label{fig:overview}
\end{figure}

% -----------------------------------------------------------------------------
